\section{Prime-shift congruences for three m\'enage sequences}
\label{sec:menage}

The entries A000033, A000159, and A000181 record three congruence conjectures
of Mark van Hoeij, dated June 10, 2019 \cite{A000033,A000159,A000181}.
We prove all three from a single polynomial identity and formalize the exact
source-definition bridge in Lean. Historical novelty is not established:
Carlitz's and Stevens's earlier papers concern m\'enage congruences, but their
complete theorem statements were unavailable for this review
\cite{Carlitz1952,Carlitz1954,Stevens1981}.

Following A058087, put $U_0(t)=1$ and, for $n\geq1$, define the classical
m\'enage hit polynomial
\begin{equation}\label{eq:menage-source}
 U_n(t)=\sum_{k=0}^{n}\frac{2n}{2n-k}
       \binom{2n-k}{k}(n-k)!(t-1)^k.
\end{equation}
Write $h_d(n)=[t^d]U_n(t)$. The triangle in A058087 lists coefficients in
descending degree, so $h_d(n)=\mathrm{A058087}(n,n-d)$ when $d\leq n$
\cite{A058087}. The relevant fixed degrees are $2,3,4$.

\begin{theorem}\label{thm:menage-shift}
For every prime $p$ and every $n\geq1$,
\begin{equation}\label{eq:menage-shift}
 U_{n+2p}(t)-2(t^p-1)U_{n+p}(t)+(t^p-1)^2U_n(t)
 \equiv0\pmod p
\end{equation}
coefficientwise. Consequently, for every $d<p$,
\[
 p\mid h_d(n)+2h_d(n+p)+h_d(n+2p).
\]
For $p=2$, this divisibility also holds for every $d<4$.
\end{theorem}

\begin{proof}
For $n\geq1$ and $j\geq0$, set
\[
 R(n,j)=\binom{n+j}{2j}+\binom{n+j-1}{2j},
\]
with a binomial coefficient equal to zero when its lower argument exceeds its
upper argument. If $j\leq n$, the identity
\[
 (n+j)\binom{n+j-1}{2j}=(n-j)\binom{n+j}{2j}
\]
shows that $R(n,j)=2n\binom{n+j}{2j}/(n+j)$. Reindexing
\eqref{eq:menage-source} by $j=n-k$ and using binomial symmetry gives the
integral polynomial identity
\begin{equation}\label{eq:menage-integral}
 U_n(t)=\sum_{j=0}^{n}j!R(n,j)(t-1)^{n-j}.
\end{equation}
In particular, all coefficients in the source polynomial are integers.

Work now in $\mathbb F_p[x]$. Frobenius implies
\[
\begin{aligned}
 &(1+x)^{a+2p}-2(1+x)^{a+p}+(1+x)^a\\
 &\qquad=(1+x)^a\bigl((1+x)^p-1\bigr)^2
 =x^{2p}(1+x)^a
\end{aligned}
\]
for every $a\geq0$. Therefore, for $r<2p$,
\begin{equation}\label{eq:menage-binomial}
 \binom{a+2p}{r}-2\binom{a+p}{r}+\binom ar\equiv0\pmod p.
\end{equation}
For $j<p$, apply \eqref{eq:menage-binomial} with $r=2j$, first with $a=n+j$
and then with $a=n+j-1$. Both upper arguments are nonnegative. Adding gives
\begin{equation}\label{eq:menage-rook}
 R(n+2p,j)-2R(n+p,j)+R(n,j)\equiv0\pmod p.
\end{equation}

Terms of \eqref{eq:menage-integral} with $j\geq p$ vanish modulo $p$ because
$p\mid j!$. Also $R(n,j)=0$ for $j>n$. Thus the sums for $n,n+p,n+2p$ may
all be compared over $0\leq j<p$, even when $n<p$. Multiply
\eqref{eq:menage-rook} by $j!(t-1)^{n+2p-j}$ and sum over that range.
The common exponent is nonnegative, and the zero terms with $j>n$ justify
completing the shorter sum. We obtain
\[
 U_{n+2p}(t)-2(t-1)^pU_{n+p}(t)+(t-1)^{2p}U_n(t)=0
 \quad\text{in }\mathbb F_p[t].
\]
Since $(t-1)^p=t^p-1$ in this ring, this is \eqref{eq:menage-shift}.
For $d<p$, taking coefficients of $t^d$ removes every term containing $t^p$
and yields the stated divisibility.

When $p=2$, the middle term is zero and $(t-1)^4=t^4+1$.
Hence $U_{n+4}+(t^4+1)U_n=0$ over $\mathbb F_2$, which gives the same
conclusion in every degree $d<4$.
\end{proof}

\begin{corollary}
Each of the following three source assertions holds on its entire domain:
\[
\begin{array}{c|c|c|c}
\text{Entry}&a(n)&\text{Index domain}&\text{Allowed primes}\\\hline
\mathrm{A000033}&h_2(n)&n\geq1&\text{all primes }p\\
\mathrm{A000159}&h_3(n)&n\geq3&p\ne3\\
\mathrm{A000181}&h_4(n)&n\geq4&p>3
\end{array}
\]
In each case, $p\mid a(n)+2a(n+p)+a(n+2p)$.
\end{corollary}

\begin{proof}
For $d=2$, every odd prime has $d<p$; for $d=3$, every prime other than
$2,3$ has $d<p$; for $d=4$, every prime greater than $3$ has $d<p$.
The remaining prime $2$ in the first two assertions is covered by the
additional $d<4$ conclusion. At $n=1$, A000033's explicit degree-two
coefficient formula gives zero, as required by its source offset.
\end{proof}

The Lean package proves equality between the literal rational rook sum
\eqref{eq:menage-source} and the integral definition
\eqref{eq:menage-integral}, including $U_0=1$. Its final declarations
\texttt{A000033\_prime\_shift}, \texttt{A000159\_prime\_shift}, and
\texttt{A000181\_prime\_shift} assert integer divisibility with exactly the
domains above. No conjectural recurrence or finite verification hypothesis is
used. These are three recorded assertions proved and formalized; the available
source review does not establish that they were historically open.
